(1) 质点位置为 \(x=\ell\sin\theta,y=-\ell\cos\theta\)。利用已指定的 \(\ell (t)\)，动能为
\[
K=\frac m2(\dot\ell^2+\ell^2\dot\theta^2),\qquad U=-mg\ell\cos\theta .
\]
拉格朗日函数 \(L=K-U\)，对 \(\theta \) 的 Euler–Lagrange 方程为
\[
\frac d{dt}(m\ell^2\dot\theta)+mg\ell\sin\theta=0,
\quad\Longrightarrow\quad
\ddot\theta+2\frac{\dot\ell}{\ell}\dot\theta+\frac g\ell\sin\theta=0.
\]
这里 \(\dot\ell=-u\)，故角速度前的系数为负，不能当作物理阻尼。

(2) 在 |\(\theta|\ll1\) 下令 \(x=\ell\theta\)。从 \(\theta=x/\ell\) 求导可得 \(\ddot x=\ell\ddot\theta+2\dot\ell\dot\theta+\ddot\ell\theta\)，因此
\[
\ddot x+\left(\frac g\ell-\frac{\ddot\ell}{\ell}\right)x=0.
\]
恒速收绳意味着 \(\ddot\ell=0\)，故 \(\ddot x+[g/(\ell_0-ut)]x=0\)。初值满足 \(x_0=\ell_0\theta_0,\dot x_0=\ell_0\dot\theta_0-u\theta_0\)。这个变频方程未必有简单三角函数解，不能假定频率恒定。

(3) 全部质点机械能为
\[
E=\frac m2(\dot\ell^2+\ell^2\dot\theta^2)-mg\ell\cos\theta .
\]
对时间求导并用角度方程，得到
\[
\dot E=m\dot\ell\ddot\ell-m\ell\dot\ell\dot\theta^2-mg\dot\ell\cos\theta
=[m(\ddot\ell-\ell\dot\theta^2)-mg\cos\theta]\dot\ell .
\]
括号是质点受到的沿远离支点方向的绳拉力代数分量 \(F_r\)，故 \(\dot E=F_r\dot\ell\)；收绳装置通常对质点做非零功，不能声称实验室机械能守恒。

(4) 令 g=0，角度方程积分为 \(\ell^2\dot\theta=\ell_0^2\dot\theta_0\)，所以
\[
\dot\theta(t)=\dot\theta_0[\ell_0/(\ell_0-ut)]^2,\quad
L_z=m\ell^2\dot\theta=m\ell_0^2\dot\theta_0.
\]
由于绳张力纯径向，无外力矩，角动量恒定；直接代回 \(\theta \) 方程两项抵消。这验证收绳系数必须是 \(2\dot\ell/\ell\) 而不是 \(\dot\ell/\ell\)。
